\documentclass[11pt]{article}

\usepackage[margin=1in]{geometry}
\usepackage{amsmath,amssymb,amsthm}
\usepackage{booktabs}
\usepackage{xcolor}
\usepackage[numbers]{natbib}
\usepackage[colorlinks=true,linkcolor=blue!60!black,citecolor=blue!60!black,urlcolor=blue!60!black]{hyperref}
\usepackage{microtype}
\usepackage{enumitem}

\emergencystretch=3em

\newcommand{\E}{\mathbb{E}}
\newcommand{\R}{\mathbb{R}}
\newcommand{\one}{\mathbf{1}}
\newcommand{\sg}{\operatorname{sg}}
\newcommand{\norm}[1]{\left\lVert #1 \right\rVert}
\newcommand{\inner}[2]{\left\langle #1,#2\right\rangle}

\newtheorem{proposition}{Proposition}
\theoremstyle{remark}
\newtheorem{remark}{Remark}

\title{First-Passage Forward--Backward Learning:\\A Bilinear Reachability Critic for Offline Goal-Conditioned Manipulation}
\author{Factored-FB project}
\date{July 22, 2026}

\begin{document}

\maketitle

\begin{abstract}

Standard Forward--Backward representation learning factorizes the discounted successor occupancy of a family of policies.

On OGBench \texttt{cube-single-play}, our occupancy-based FB variants reach only approximately $0.10$ mean success.

We instead retain the bilinear FB parameterization but replace the successor-occupancy objective with an absorbing first-passage Bellman objective.

The resulting critic estimates the expected discount to the first time that the cube enters the goal region.

Unlike the continuous-state successor-measure objective, this formulation provides a directly computable binary TD target for every state--goal pair and produces values bounded in $[0,1]$.

With the FB covariance regularizer disabled, the method reaches $0.967$ mean success over three seeds ($0.956$, $0.964$, and $0.980$), compared with approximately $0.94$ for the bilinear GCIQL baseline.

The method is best interpreted as an FB-factorized goal-conditioned reachability critic rather than as an estimator of the canonical FB successor measure.

\end{abstract}

\section{Setting}

We consider an offline goal-conditioned Markov decision process $(\mathcal S,\mathcal A,P,\gamma)$.

The offline dataset contains transitions $(s,a,s')\sim\rho$ without task rewards.

A goal $g$ is represented by a target observation.

Let $c(s)\in\R^3$ denote the cube position contained in observation $s$.

In the implementation, $c(s)$ corresponds to \texttt{obs[19:22]}.

We use the bilinear critic

\begin{equation}
Q_\theta(s,a,g)=\frac{1}{\sqrt d}\inner{F_\theta(s,a)}{B_\theta(g)},
\label{eq:bilinear}
\end{equation}

where $F_\theta(s,a)\in\R^d$ is the forward representation and $B_\theta(g)\in\R^d$ is the goal representation.

The experiments use $d=50$ and two critic heads.

The actor is trained to maximize this critic while remaining close to actions in the offline dataset.

\section{Relation to standard Forward--Backward representations}

Forward--Backward representation learning factorizes the discounted successor measure of a policy \citep{touati2021learning}.

For a fixed policy $\pi$, the state-goal successor occupancy may be written informally as

\begin{equation}
M_{\mathrm{occ}}^\pi(s,a;g)=\E_\pi\left[\sum_{t\geq 0}\gamma^t\delta_g(s_{t+1})\mid s_0=s,a_0=a\right].
\label{eq:occupancy}
\end{equation}

FB learns forward and backward representations whose inner product approximates the density of this measure relative to the data distribution.

This factorization permits a task vector to be constructed after representation learning and used to extract a policy.

The standard successor-occupancy objective is not identical to the episodic goal-reaching objective used by OGBench.

Occupancy counts every discounted visit to the goal, whereas OGBench evaluates whether the object reaches the goal region.

For a unit indicator, discounted occupancy is bounded by $1/(1-\gamma)$.

Nevertheless, its density relative to the offline data distribution can be large in low-density regions, and the factorization remains non-identifiable under reciprocal rescaling of $F$ and $B$.

In continuous state spaces, exact equality with a sampled goal also does not provide an ordinary binary regression target.

Canonical FB therefore uses a successor-density objective with a special diagonal term rather than standard regression onto an observed sparse reward.

\section{First-passage reachability}

\subsection{Goal region}

We define a goal region using the cube position.

\begin{equation}
\mathcal G(g)=\left\{s'\in\mathcal S:\norm{c(s')-c(g)}_2\leq\varepsilon\right\}.
\label{eq:goal-set}
\end{equation}

The corresponding goal-reaching reward is

\begin{equation}
r_g(s')=\one\left[\norm{c(s')-c(g)}_2\leq\varepsilon\right].
\label{eq:hit-reward}
\end{equation}

The implementation uses $\varepsilon=0.4$ in the scaled observation space.

This corresponds to the $0.04$ metre success threshold used by the OGBench Cube environment \citep{park2025ogbench}.

\subsection{Absorbing Bellman operator}

We terminate the value recursion when the successor state enters the goal region.

For a goal-conditioned policy $\pi(a\mid s,g)$, define

\begin{equation}
(\mathcal T^\pi Q)(s,a,g)=\E_{s'\sim P(\cdot\mid s,a),\,a'\sim\pi(\cdot\mid s',g)}\left[r_g(s')+\gamma\left(1-r_g(s')\right)Q(s',a',g)\right].
\label{eq:operator}
\end{equation}

When $r_g(s')=1$, the target is exactly $1$ and no future value is added.

When $r_g(s')=0$, the target is the discounted value of continuing from $s'$.

\begin{proposition}
The operator $\mathcal T^\pi$ is a $\gamma$-contraction under the supremum norm.
\end{proposition}

\begin{proof}

For any $Q_1$ and $Q_2$,

\begin{equation}
\left|(\mathcal T^\pi Q_1)(s,a,g)-(\mathcal T^\pi Q_2)(s,a,g)\right|\leq\gamma\E\left[(1-r_g(s'))\left|Q_1(s',a',g)-Q_2(s',a',g)\right|\right].
\end{equation}

Since $1-r_g(s')\in\{0,1\}$,

\begin{equation}
\norm{\mathcal T^\pi Q_1-\mathcal T^\pi Q_2}_\infty\leq\gamma\norm{Q_1-Q_2}_\infty.
\end{equation}

\end{proof}

\subsection{First-passage interpretation}

Let

\begin{equation}
T_g=\inf\left\{t\geq 1:s_t\in\mathcal G(g)\right\}
\label{eq:hitting-time}
\end{equation}

denote the first time that the trajectory enters the goal region.

We set $T_g=\infty$ if the goal is never reached.

\begin{proposition}
The unique fixed point of \eqref{eq:operator} is

\begin{equation}
Q_{\mathrm{FP}}^\pi(s,a,g)=\E_\pi\left[\gamma^{T_g-1}\one[T_g<\infty]\mid s_0=s,a_0=a\right].
\label{eq:first-passage}
\end{equation}

\end{proposition}

\begin{proof}

Unrolling the absorbing Bellman recursion gives

\begin{equation}
Q_{\mathrm{FP}}^\pi(s,a,g)=\E_\pi\left[\sum_{t\geq 1}\gamma^{t-1}\left(\prod_{u=1}^{t-1}(1-r_g(s_u))\right)r_g(s_t)\right].
\end{equation}

The product is one only while the goal has not previously been reached.

Consequently, exactly one term is nonzero, at $t=T_g$.

This yields \eqref{eq:first-passage}.

\end{proof}

The first-passage value belongs to $[0,1]$.

It is one for immediate success, decreases with the time required to reach the goal, and is zero when the goal is never reached.

It is therefore directly aligned with goal-reaching control.

\section{Batched training objective}

Consider a minibatch

\begin{equation}
\left\{(s_i,a_i,s_i')\right\}_{i=1}^N.
\end{equation}

We use the next observations in the batch as candidate goals,

\begin{equation}
g_j=s_j'.
\end{equation}

This produces an $N\times N$ matrix of state--action and goal pairs.

\subsection{Critic matrix}

The online critic matrix is

\begin{equation}
Q_{ij}=\frac{1}{\sqrt d}\inner{F_\theta(s_i,a_i)}{B_\theta(g_j)}.
\label{eq:q-matrix}
\end{equation}

For the exact goal-conditioned Bellman operator, the target action for entry $(i,j)$ should be sampled from the policy conditioned on goal $g_j$,

\begin{equation}
a_{ij}^+\sim\pi_{\bar\theta}(\cdot\mid s_i',g_j).
\label{eq:pair-action}
\end{equation}

The corresponding target critic is

\begin{equation}
\bar Q_{ij}=\sg\left[\frac{1}{\sqrt d}\inner{F_{\bar\theta}(s_i',a_{ij}^+)}{B_{\bar\theta}(g_j)}\right].
\label{eq:target-q}
\end{equation}

\begin{remark}
The current vectorized implementation appears to construct one target action $a_i^+$ per row and reuse it across all goal columns.
If so, this is a shared-action approximation to \eqref{eq:pair-action}.
It is exact on the diagonal when the row action is conditioned on $g_i$, but it is not the exact goal-conditioned target for arbitrary off-diagonal goals.
\end{remark}

\subsection{All-pairs goal labels}

The goal-reaching matrix is

\begin{equation}
H_{ij}=\one\left[\norm{c(s_i')-c(g_j)}_2\leq\varepsilon\right].
\label{eq:hit-matrix}
\end{equation}

Since $g_i=s_i'$, every diagonal entry satisfies $H_{ii}=1$.

Off-diagonal entries may also be positive when two next observations contain nearby cube positions.

The Bellman target is

\begin{equation}
Y_{ij}=H_{ij}+\gamma(1-H_{ij})\bar Q_{ij}.
\label{eq:target}
\end{equation}

The TD residual is

\begin{equation}
\Delta_{ij}=Q_{ij}-Y_{ij}.
\label{eq:residual}
\end{equation}

\subsection{Diagonal and off-diagonal losses}

Let $I$ denote the identity matrix and let $\mathrm{Off}=\mathbf 1-I$.

The off-diagonal loss is

\begin{equation}
\mathcal L_{\mathrm{off}}=\frac{0.5}{\sum_{ij}\mathrm{Off}_{ij}}\sum_{i\neq j}\Delta_{ij}^2.
\label{eq:off-loss}
\end{equation}

The diagonal loss is

\begin{equation}
\mathcal L_{\mathrm{diag}}=\frac{1}{N}\sum_i\Delta_{ii}^2.
\label{eq:diag-loss}
\end{equation}

Unlike the canonical FB diagonal attraction term, the first-passage diagonal is trained by squared TD regression.

The diagonal target is exactly one because each transition reaches its own next-state goal in one step.

\subsection{Loss normalization}

The implementation normalizes both TD terms using the stop-gradient critic scale

\begin{equation}
S_Q=\sg\left[\frac{1}{N^2}\sum_{ij}Q_{ij}^2\right]+\epsilon.
\label{eq:q-scale}
\end{equation}

The normalized losses are

\begin{equation}
\widetilde{\mathcal L}_{\mathrm{off}}=\frac{\mathcal L_{\mathrm{off}}}{S_Q},
\qquad
\widetilde{\mathcal L}_{\mathrm{diag}}=\frac{\mathcal L_{\mathrm{diag}}}{S_Q}.
\label{eq:normalized-losses}
\end{equation}

This normalization is an optimization heuristic and does not change the first-passage fixed point.

\section{Backward-representation regularization}

The implementation optionally applies an FB-style covariance regularizer to the backward representations.

Let $\Psi\in\R^{N\times d}$ contain the goal representations in a minibatch and define

\begin{equation}
C=\Psi\Psi^\top.
\end{equation}

The implemented regularizer is

\begin{equation}
\mathcal L_{\mathrm{cov}}=\frac{0.5}{\sum_{ij}\mathrm{Off}_{ij}}\sum_{i\neq j}C_{ij}^2-\frac{1}{N}\sum_i C_{ii}.
\label{eq:cov-loss}
\end{equation}

The total critic loss is

\begin{equation}
\mathcal L_{\mathrm{critic}}=\widetilde{\mathcal L}_{\mathrm{off}}+\widetilde{\mathcal L}_{\mathrm{diag}}+\lambda_{\mathrm{cov}}\mathcal L_{\mathrm{cov}}.
\label{eq:critic-loss}
\end{equation}

This term is referred to as orthonormalization in the implementation, although it is more precisely a covariance-shaping regularizer.

It discourages correlations between different goal embeddings and encourages nonzero embedding norms.

\section{Policy learning}

For a goal observation $g$, the task representation is

\begin{equation}
z_g=\sg\left[\frac{1}{P}\sum_{p=1}^P B_{\theta,p}(g)\right],
\label{eq:task-vector}
\end{equation}

where $P$ is the number of critic heads.

The actor uses the bilinear value

\begin{equation}
Q_\theta(s,a\mid g)=\frac{1}{\sqrt d}\inner{F_\theta(s,a)}{z_g}.
\label{eq:actor-q}
\end{equation}

A flow or DDPG+BC actor maximizes this value while remaining close to the offline action distribution.

Because \eqref{eq:first-passage} decreases with the first-passage time, maximizing \eqref{eq:actor-q} favors actions that reach the goal region quickly.

\section{Why the method works on \texttt{cube-single-play}}

\subsection{Alignment with the evaluation objective}

OGBench Cube tasks declare success when the cube position is within $0.04$ metres of the target position \citep{park2025ogbench}.

The goal-conditioned environment terminates after success and reports a binary success signal.

The first-passage reward uses the same position coordinates and the same tolerance.

The training target therefore closely matches the evaluation criterion.

Standard occupancy learning does not encode this terminal first-arrival semantics.

\subsection{Low-dimensional goal structure}

The complete observation contains the robot state, gripper state, object state, and other dynamical information.

However, success in \texttt{cube-single} depends only on the three-dimensional cube position.

Many full goal observations should therefore correspond to the same goal-reaching task.

The useful goal space is a quotient of the full observation space under the equivalence relation

\begin{equation}
g_1\sim g_2\quad\Longleftrightarrow\quad c(g_1)\approx c(g_2).
\label{eq:goal-equivalence}
\end{equation}

Without covariance regularization, $B$ can learn to ignore goal-observation variables that do not affect the cube-position reward.

This is consistent with the FB observation that the backward map may operate only on reward-relevant state features when those features are known \citep{touati2021learning}.

\subsection{A low-rank reachability function}

The task goal is parameterized by a three-dimensional position, while the latent dimension is $d=50$.

The first-passage value also varies smoothly over nearby cube positions except at the success boundary.

Its state--action--goal value matrix can therefore have a much lower effective rank than the nominal latent dimension.

A free bilinear factorization can allocate most of its variation to a small correlated subspace.

This is sufficient to represent spatial reachability while allowing $F$ to retain the robot and contact information needed to select actions.

\subsection{All-pairs hindsight supervision}

Each minibatch produces $N^2$ state--goal pairs.

The diagonal provides $N$ guaranteed one-step positive examples.

Nearby off-diagonal pairs provide additional positive examples.

The remaining pairs train discounted reachability through bootstrapping.

This gives the critic substantially more structured supervision than an exact-state successor-density objective in continuous space.

\subsection{Behavior support}

The OGBench play datasets contain diverse manipulation behavior generated by scripted data-collection policies \citep{park2025ogbench}.

The dataset includes reaching, grasping, transporting, and releasing the cube at many positions.

The behavior-constrained actor can therefore construct successful policies largely from actions already represented in the dataset.

The critic's main task is to rank these supported actions according to first-passage reachability.

\section{Why covariance regularization hurts}

The covariance regularizer is useful in canonical FB partly because the factorization is non-identifiable and task vectors may be sampled independently from an isotropic latent distribution.

Neither motivation applies directly to the present objective.

Tasks are represented by $z_g=B(g)$ rather than by independent isotropic vectors.

The actor depends only on the product $F(s,a)^\top B(g)$.

The first-passage TD target directly constrains this product.

Nearby goals should also have similar value functions and therefore similar embeddings.

Forcing these embeddings toward decorrelation conflicts with the smooth geometry of the cube-position goal space.

The conflict is stronger because $B$ receives a full goal observation even though the reward depends only on three position coordinates.

Without covariance regularization, $B$ may collapse nuisance differences in robot configuration, gripper state, velocity, and contact state.

With covariance regularization, the representation is encouraged to distinguish goal observations that should be equivalent under the task reward.

The experimental results are consistent with this explanation:

\begin{center}
\begin{tabular}{lc}
\toprule
Method & Mean success \\
\midrule
Occupancy FB variants & $\approx 0.10$ \\
First-passage with covariance regularization & $0.225$ \\
First-passage without covariance regularization & $0.967$ \\
Bilinear GCIQL & $\approx 0.94$ \\
\bottomrule
\end{tabular}
\end{center}

Removing the covariance regularizer alone does not rescue the occupancy objective.

The first-passage target supplies the correct goal-reaching semantics, while removing the regularizer allows the backward representation to use the corresponding low-rank and nuisance-invariant structure.

\section{Interpretation and limitations}

The method retains three main components of FB:

\begin{enumerate}[leftmargin=1.5em]
\item a forward representation of state--action pairs;
\item a backward representation of goals;
\item a bilinear value used for task-conditioned policy extraction.
\end{enumerate}

It does not retain the canonical FB successor-measure objective.

Its fixed point is a goal-conditioned first-passage value function.

The method should therefore be described as \emph{first-passage FB}, \emph{an FB-factorized reachability critic}, or \emph{bilinear first-passage Q-learning}.

It should not be claimed to preserve the standard FB guarantee of recovering optimal policies for arbitrary rewards.

The current construction also uses task-specific information: the cube-position coordinates and the OGBench success threshold.

This specialization is a likely source of its strong performance.

The following diagnostics would directly test the proposed explanation:

\begin{enumerate}[leftmargin=1.5em]
\item measure the singular-value spectrum of the learned state--goal value matrix;
\item measure the covariance spectrum of $B(g)$ with and without regularization;
\item train $B$ using only the three cube-position coordinates;
\item vary the latent dimension over $d\in\{3,5,10,20,50\}$;
\item restore covariance regularization after restricting $B$ to cube position;
\item compare pair-specific and row-shared bootstrap actions;
\item evaluate the same objective on multi-object and Scene tasks.
\end{enumerate}

A particularly informative prediction is that a position-only backward map should substantially reduce the gap between regularized and unregularized training.

\section{Code correspondence}

All line references below refer to \texttt{tools/wandb\_mode\_shim/gciql\_flow.py}.

\begin{center}
\footnotesize
\setlength{\tabcolsep}{4pt}
\begin{tabular}{@{}lll@{}}
\toprule
Mathematical object & Implementation & Lines \\
\midrule
$Q_{ij}$ & \texttt{einsum('hid,hjd->hij', phi, psi) / sqd} & 1157 \\
$\bar Q_{ij}$ & target forward--backward inner product & 1176 \\
$H_{ij}$ & \texttt{dist <= 0.4} on \texttt{nobs[:,19:22]} & 1200--1203 \\
$\Delta_{ij}$ & \texttt{M - hit - gamma * (1-hit) * tM} & 1090 \\
$\mathcal L_{\mathrm{diag}}$ & \texttt{mean(diag(diff)**2)} & 1091 \\
$\mathcal L_{\mathrm{off}}$ & \texttt{0.5 * sum((diff * off)**2) / sum(off)} & 1092 \\
$S_Q$ & \texttt{sg(mean(M**2))} & 1093--1097 \\
$\mathcal L_{\mathrm{cov}}$ & covariance off-diagonal and diagonal terms & 1315--1318 \\
$\mathcal L_{\mathrm{critic}}$ & TD terms plus covariance regularization & 1320 \\
Total loss & \texttt{measure\_coef * m\_loss + actor\_loss} & 1677 \\
$z_g$ & \texttt{sg(psi.mean(axis=0))} & 425, 467 \\
\bottomrule
\end{tabular}
\end{center}

\section{Experimental configuration}

The successful configuration uses

\begin{itemize}[leftmargin=1.5em]
\item \texttt{value\_arch=bilinear};
\item \texttt{value\_latent\_dim=50};
\item \texttt{measure\_only=true};
\item \texttt{measure\_coef=1.0};
\item \texttt{measure\_legs=critic};
\item \texttt{measure\_first\_passage=true};
\item \texttt{measure\_scale\_norm=true};
\item \texttt{measure\_ortho\_coef=0.0};
\item \texttt{bc\_coeff=3.0};
\item a flow actor;
\item $500\,000$ gradient steps.
\end{itemize}

The corresponding experiment group is

\begin{center}
\texttt{GCIQL-State-Flow-Measure-FirstPassage-NoOrtho-d50}.
\end{center}

\section{Conclusion}

The successful modification does not merely stabilize the original FB loss.

It changes the learned object from discounted successor occupancy to discounted first-arrival reachability.

This produces a bounded Bellman target that matches the terminal goal-reaching objective used by OGBench Cube.

The large improvement without covariance regularization is consistent with the low-dimensional and highly correlated structure of cube-position goals.

The result suggests that, for manipulation, an FB-style bilinear architecture can be effective when its training objective and goal representation are aligned with the task-relevant notion of success.

\bibliographystyle{plainnat}

\begin{thebibliography}{9}

\bibitem[Touati and Ollivier(2021)]{touati2021learning}
Ahmed Touati and Yann Ollivier.
\newblock Learning one representation to optimize all rewards.
\newblock In \emph{Advances in Neural Information Processing Systems}, 2021.

\bibitem[Park et~al.(2025)]{park2025ogbench}
Seohong Park, Kevin Frans, Benjamin Eysenbach, and Sergey Levine.
\newblock OGBench: Benchmarking offline goal-conditioned reinforcement learning.
\newblock In \emph{International Conference on Learning Representations}, 2025.

\end{thebibliography}

\end{document}